\documentclass[11pt,a4paper]{article}

\usepackage[utf8]{inputenc}
\usepackage[T1]{fontenc}
\usepackage{amsmath,amssymb,amsthm}
\usepackage{bm}
\usepackage{mathtools}
\usepackage{geometry}
\geometry{margin=1in}
\usepackage{hyperref}
\usepackage{enumitem}

\DeclareMathOperator{\SO}{SO}
\DeclareMathOperator{\SE}{SE}
\DeclareMathOperator{\softmax}{softmax}
\DeclareMathOperator{\softplus}{softplus}
\DeclareMathOperator{\diag}{diag}
\DeclareMathOperator{\Span}{span}
\DeclareMathOperator{\sgn}{sgn}
\DeclareMathOperator{\SiLU}{SiLU}
\DeclareMathOperator{\Sigmoid}{Sigmoid}
\DeclareMathOperator{\SwiGLU}{SwiGLU}
\DeclareMathOperator{\GELU}{GELU}
\DeclareMathOperator{\LN}{LayerNorm}
\DeclareMathOperator{\RMSN}{RMSNorm}
\DeclareMathOperator{\Embed}{Embed}
\DeclareMathOperator{\softcap}{softcap}
\DeclareMathOperator*{\SUM}{\textstyle\sum}

\newcommand{\R}{\mathbb{R}}
\newcommand{\N}{\mathbb{N}}
\newcommand{\Z}{\mathbb{Z}}
\newcommand{\C}{\mathbb{C}}
\newcommand{\bA}{\bm{A}}
\newcommand{\bU}{\bm{U}}
\newcommand{\bV}{\bm{V}}
\newcommand{\bG}{\bm{G}}
\newcommand{\Sph}{S^2}
\newcommand{\bR}{\bm{R}}
\newcommand{\br}{\bm{r}}
\newcommand{\brh}{\hat{\bm{r}}}
\newcommand{\bx}{\bm{x}}
\newcommand{\by}{\bm{y}}
\newcommand{\bz}{\bm{z}}
\newcommand{\bmF}{\bm{F}}
\newcommand{\bh}{\bm{h}}
\newcommand{\bq}{\bm{q}}
\newcommand{\bk}{\bm{k}}
\newcommand{\bv}{\bm{v}}
\newcommand{\bg}{\bm{g}}
\newcommand{\bu}{\bm{u}}
\newcommand{\Te}{\bm{T}}
\newcommand{\bD}{\bm{D}}
\newcommand{\bW}{\bm{W}}
\newcommand{\bP}{\bm{P}}
\newcommand{\bphi}{\bm{\phi}}
\newcommand{\brho}{\bm{\rho}}
\newcommand{\balpha}{\bm{\alpha}}
\newcommand{\bxi}{\bm{\xi}}
\newcommand{\rcut}{r_{\mathrm{c}}}
\newcommand{\lmax}{l_{\max}}
\newcommand{\mmax}{m_{\max}}
\newcommand{\Nb}{\mathcal{N}}
\newcommand{\Hb}{\mathcal{H}}

\theoremstyle{plain}
\newtheorem{theorem}{Theorem}
\newtheorem{proposition}[theorem]{Proposition}
\newtheorem{lemma}[theorem]{Lemma}
\theoremstyle{definition}
\newtheorem{definition}[theorem]{Definition}
\theoremstyle{remark}
\newtheorem{remark}[theorem]{Remark}

\title{The EquiformerV3 Architecture:\\ A Mathematical Construction\\ {\large With a Component-by-Component Comparison to DPA4}}
\author{}
\date{}

\begin{document}
\maketitle

\begin{abstract}
We give a mathematically explicit construction of the EquiformerV3 graph neural network as a parametric family of $\SO(3)$-equivariant functions on atomic configurations.
We follow the source distribution at \texttt{experimental/models/equiformer\_v3/} of the EquiformerV3 release, transcribing the implementation into definitions, operators, and named blocks rather than into code.
The construction proceeds through the configuration space, edge graph, edge-aligned $\SO(3)$ frame, Wigner-$D$ rotation to and from that frame, real spherical-harmonic feature layout, $\SO(2)$ linear operators in the local frame, equivariant attention, $\Sph$-grid feed-forward network, equivariant layer normalisation, pre-norm transformer block, and energy/force/stress output heads.
The final section places EquiformerV3 side-by-side with DPA4 of the companion document and identifies, equation by equation, the choices on which the two architectures agree and the choices on which they diverge.
\end{abstract}

\tableofcontents

% =============================================================
\section{Overview of the Architecture}\label{sec:ov}

EquiformerV3 is a parametric family of $\SO(3)$-equivariant functions
\begin{equation}\label{eq:overall}
\bigl(Z,R\bigr)\;\longmapsto\;\bigl(E_\Theta(Z,R),\,\bmF_\Theta(Z,R),\,\sigma_\Theta(Z,R)\bigr)
\end{equation}
of atomic configurations $(Z,R)$ with species $Z\in\{1,\dots,Z_{\max}\}^N$ and positions $R\in\R^{3N}$.
The energy $E_\Theta(Z,R)\in\R$ is a translation-, permutation-, and rotation-invariant scalar; the forces $\bmF_\Theta(Z,R)\in\R^{3N}$ are $\SO(3)$-equivariant vectors; the stress $\sigma_\Theta(Z,R)\in\R^{3\times 3}$ is an $\SO(3)$-equivariant symmetric tensor.
Forces and stress can be obtained either by autograd through $E_\Theta$ (\emph{gradient prediction}) or by separate $\SO(3)$-equivariant output heads (\emph{direct prediction}); the choice is a configuration flag.

\paragraph{High-level pipeline.}
Given an atomic configuration $(Z,R)$, EquiformerV3 executes the following sequence:
\begin{enumerate}[leftmargin=2em]
\item Build the directed edge graph $\Nb=\{(i,j):0<r_{ij}<\rcut\}$, possibly with periodic images, capped at $\Nb_{\max}$ neighbours per atom.
\item For each edge $(i,j)$ construct a $3\times 3$ rotation matrix $\bR_{ij}\in\SO(3)$ whose middle column is the unit edge vector $\brh_{ij}$; this defines the edge-aligned local frame in which $\brh_{ij}$ becomes the local $\by$-axis.
\item For each edge compute the block-diagonal Wigner-$D$ matrix $\bD_{ij}\in\R^{D\times D}$, $D=(\lmax+1)^2$, in a truncated $m$-primary layout (\S\ref{sec:so3rot}).
\item Initialise per-atom features $\bh_i\in\R^{D\times C}$ from the atomic-number embedding (only the $l=0$ slot is non-zero), then add a one-shot \emph{edge-degree embedding} that promotes the $m=0$ radial channel up to all $(l,m)$ slots via $\bD_{ij}^{-1}$.
\item Apply $L$ pre-norm transformer blocks (\S\ref{sec:block}), each consisting of an equivariant graph-attention sublayer (\S\ref{sec:attn}) and an equivariant feed-forward sublayer (\S\ref{sec:ffn}), with $\Sph$-grid nonlinearities.
\item Apply a final equivariant layer norm; the $l=0$ slot of each atom is fed through a scalar MLP to produce a per-atom energy, summed over the structure with a fixed normaliser.
\item For forces and stress: \emph{(direct)} apply an extra equivariant-attention head whose $l=1$ output gives the force vector, and an equivariant FFN head whose first $(\lmax+1)^2$ irreducible coefficients are pulled back through a fixed change-of-basis matrix to a symmetric stress tensor; or \emph{(gradient)} differentiate $E_\Theta$ with respect to positions and a symmetric displacement tensor by autograd.
\end{enumerate}

\paragraph{Two design ideas, in plain language.}
First, equivariance is enforced by carrying out all linear operations either as $\SO(3)$ scalar-block channel mixers $L^{\mathrm{ch}}$, or, after a per-edge rotation to the local frame, as $\SO(2)_{\by}$-equivariant linear operators $L^{\mathrm{SO2}}$.
Second, nonlinearities that would otherwise destroy equivariance act on a sampled $\Sph$-grid: $\SO(3)$ features are evaluated on a latitude--longitude grid via the real spherical-harmonic basis, an unconstrained pointwise MLP acts on grid samples, and the result is projected back into the irreducible basis.
No Clebsch--Gordan coefficient appears anywhere in the parameterisation: \emph{equivariant--equivariant} tensor products are never computed; only \emph{equivariant--invariant} multiplications by edge-conditioned scalar radials, and \emph{equivariant--equivariant} compositions through pointwise multiplication on the $\Sph$-grid, are allowed.

\paragraph{Notation.}
Throughout, $i,j\in\{1,\dots,N\}$ are atom indices, $(i,j)\in\Nb$ are directed edges with source $i$ and target $j$, $E=|\Nb|$ is the number of edges, $C$ is the feature width (\texttt{num\_channels}), $H$ is the number of attention heads (\texttt{num\_heads}), $\lmax$ is the maximum $\SO(3)$ degree, and $\mmax\le\lmax$ is the maximum $\SO(2)$ order retained after rotation.
Default hyperparameters in the public release are $L=12$, $C=128$, $H=8$, $\lmax=6$, $\mmax=2$, $\rcut=12$\,\AA{}, $\Nb_{\max}=20$.

% =============================================================
\section{Configuration Space and Symmetries}\label{sec:conf}

\paragraph{Atomic configuration.}
An atomic configuration is a pair $(Z,R)$ with $Z=(Z_1,\dots,Z_N)\in\{1,\dots,Z_{\max}\}^N$ and $R=(\bR_1,\dots,\bR_N)\in\R^{3N}$.
Periodic images, when present, are generated by integer combinations of the simulation-cell columns and contribute additional edge candidates.

\paragraph{Edge data.}
For an ordered pair $(i,j)$ with $j\ne i$ (modulo periodic images), define
\begin{equation}\label{eq:edgevec}
\br_{ij}\;=\;\bR_j-\bR_i,\qquad r_{ij}=\|\br_{ij}\|_2,\qquad \brh_{ij}=\br_{ij}/r_{ij}.
\end{equation}
The directed edge set is $\Nb=\{(i,j):0<r_{ij}<\rcut,\;j\in\text{top-}\Nb_{\max}(i)\}$ where the top-$\Nb_{\max}$ selection retains the closest $\Nb_{\max}$ neighbours of $i$, breaking ties deterministically when so configured (flag \texttt{enforce\_max\_neighbors\_strictly}).

\paragraph{Symmetry group and action.}
The relevant symmetry group is $\SE(3)\times S_N$ (Euclidean isometries times atom-index permutations), restricted to the rotation subgroup $\SO(3)\subset\SE(3)$ for the equivariant features (translations factor out because only $\br_{ij}$ appears; reflections are not enforced).
A rotation $g\in\SO(3)$ acts on positions by $g\cdot\bR_i\;=\;g\bR_i$, on edge vectors by $g\cdot\br_{ij}=g\br_{ij}$, and on a degree-$l$ feature $\bh^{(l,\cdot)}_i\in\R^{2l+1}$ by the real Wigner-$D$ matrix of degree $l$:
\begin{equation}\label{eq:wignerD-action}
(g\cdot\bh^{(l)}_i)_m \;=\;\sum_{m'} D^{(l)}_{m,m'}(g)\,\bh^{(l,m')}_i,
\end{equation}
the sum running over $m'\in\{-l,\dots,l\}$.
The full per-atom feature transforms by the block-diagonal Wigner matrix
\begin{equation}\label{eq:Dgblock}
\bD(g)\;=\;\mathrm{blockdiag}\bigl(D^{(0)}(g),D^{(1)}(g),\dots,D^{(\lmax)}(g)\bigr)\in\R^{D\times D},\qquad D=(\lmax+1)^2.
\end{equation}

% =============================================================
\section{Edge Frame and Wigner-$D$ Rotation}\label{sec:so3rot}

\subsection{\texorpdfstring{The edge rotation matrix $\bR_{ij}$}{The edge rotation matrix R\_ij}}\label{sec:edgerot}

EquiformerV3 places each edge into a coordinate frame in which the unit edge vector $\brh_{ij}$ becomes the local $\by$-axis.
The construction is procedural, not analytic, and uses a fresh random perpendicular per edge.
Concretely (see \texttt{edge\_rot\_mat.py}):

\begin{enumerate}[leftmargin=2em]
\item Set $\bu_x=\brh_{ij}\in\Sph$.
\item Draw $\bu^{(0)}\sim\mathrm{Uniform}([-\tfrac12,\tfrac12]^3)$, normalise to $\Sph$; form three candidates $\bu^{(0)},\bu^{(1)},\bu^{(2)}$ from $\bu^{(0)}$ by axis-aligned permutations, and select $\bu^\star=\arg\min_{k}|\bu^{(k)}\cdot\bu_x|$.
This guarantees $|\bu^\star\cdot\bu_x|<0.99$.
\item Compute $\bu_z=(\bu_x\times\bu^\star)/\|\bu_x\times\bu^\star\|$ and $\bu_y=-(\bu_x\times\bu_z)/\|\bu_x\times\bu_z\|$.
\item Assemble
\begin{equation}\label{eq:edgerot}
\bR_{ij}\;\coloneqq\;\bigl[\;\bu_z\;\big|\;\bu_x\;\big|\;\bu_y\;\bigr]^{\top}\in\SO(3),
\end{equation}
so that $\bR_{ij}\,\bR_{ij}^{\top}=I_3$ and the middle column of $\bR_{ij}^{\top}$ is $\bu_x=\brh_{ij}$.
\end{enumerate}

This choice of frame has two important consequences.
First, the local $\by$-axis carries the bond direction, so the subgroup of $\SO(3)$ stabilising the edge is the rotations about $\by$, i.e.\ $\SO(2)_{\by}$; this is the group with respect to which all in-frame linear operators must be equivariant.
Second, $\bR_{ij}$ contains the random direction $\bu^\star$ in its first and third columns; the random choice cancels in any object that is $\SO(2)_{\by}$-invariant, but only those objects.
EquiformerV3 therefore restricts the in-frame linear operators to be $\SO(2)_{\by}$-equivariant.

\begin{remark}\label{rmk:rand-frame}
The randomness in $\bu^\star$ has no observable effect provided every operator performed in the local frame is exactly $\SO(2)_{\by}$-equivariant.
The code achieves this by working in the $m$-primary layout where the $\SO(2)_{\by}$-action is diagonal in $m$ (one $2\times 2$ block per $|m|>0$, one $1\times 1$ block per $m=0$): linear maps are constrained to act independently in each $|m|$-stratum (\S\ref{sec:so2lin}); pointwise grid operations are applied uniformly across longitudes (\S\ref{sec:ffn}); attention logits are $\SO(2)$-invariant scalars by construction (\S\ref{sec:attn}).
\end{remark}

\subsection{Wigner-$D$ from a $3\times 3$ rotation matrix}\label{sec:wignerD}

Given $\bR_{ij}$, EquiformerV3 computes the block-diagonal Wigner-$D$ matrix that rotates $\SO(3)$ features by $\bR_{ij}^{-1}$ (so that the local-frame view of a feature is obtained by left-multiplication).
The procedure (see \texttt{so3.py:\_rotation\_to\_wigner\_matrix}) is:
\begin{enumerate}[leftmargin=2em]
\item Read off the unit vector $\bxi\;=\;[\bR_{ij}]_{:,1}\in\Sph$ (the second column of $\bR_{ij}$ in row-major indexing), which equals $\bu_x=\brh_{ij}$.
\item Compute the ZYZ Euler angles $(\alpha,\beta,\gamma)$ such that
\begin{equation}\label{eq:zyz-angles}
\bR_{ij}\;=\;R_z(\alpha)\,R_y(\beta)\,R_z(\gamma),
\end{equation}
using $(\alpha,\beta)=\mathrm{xyz\_to\_angles}(\bxi)$ and then $\gamma=\mathrm{atan2}\bigl([R_z(\alpha)R_y(\beta)]^{\top}\bR_{ij}\bigr)_{0,2},\,\bigl([R_z(\alpha)R_y(\beta)]^{\top}\bR_{ij}\bigr)_{0,0}\bigr)$.
\item For each $l\in\{0,\dots,\lmax\}$ assemble the degree-$l$ real Wigner-$D$ block
\begin{equation}\label{eq:wignerD-real}
D^{(l)}_{m,m'}(\alpha,\beta,\gamma)\;=\;\bigl[e^{-i\alpha\,M_z^{(l)}}\,d^{(l)}(\beta)\,e^{-i\gamma\,M_z^{(l)}}\bigr]_{m,m'}^{\mathrm{re}},
\end{equation}
where $d^{(l)}(\beta)\in\R^{(2l+1)\times(2l+1)}$ is the reduced Wigner-$d$ for the small rotation $R_y(\beta)$, contracted against the precomputed Jucys--Murphy tensor $J_d$ stored in \texttt{Jd.pt}, and the real-basis representation eliminates complex arithmetic.
\item Concatenate into $\bD_{ij}=\mathrm{blockdiag}(D^{(0)},D^{(1)},\dots,D^{(\lmax)})\in\R^{D\times D}$.
\end{enumerate}

\subsection{$m$-primary truncated layout}\label{sec:mprim}

For $\mmax<\lmax$ the in-frame computations retain only those components with $|m|\le\mmax$.
Let
\begin{equation}\label{eq:Dm}
D_m\;\coloneqq\;\sum_{l=0}^{\lmax}\bigl(2\min(l,\mmax)+1\bigr)
\end{equation}
be the size of the truncated layout.
The $m$-primary mapping $\bm{T}_m\in\R^{D_m\times D}$ permutes indices so that the truncated stream
$(0,0,0,\dots,0;\;1^+_l,\dots,1^-_l;\;\dots;\;m^+,\dots,m^-)$
groups all entries with the same $|m|$ together (real parts in front of imaginary parts, in $l$-increasing order within each block).
In this layout the $\SO(2)_{\by}$-action is block-diagonal: a $1\times 1$ block on the $m=0$ stratum and a $2\times 2$ rotation block on each $|m|>0$ stratum.

The truncated Wigner-$D$ in $m$-primary layout, used to rotate the global-frame feature into the local frame, is
\begin{equation}\label{eq:wignerD-trunc}
\widetilde\bD_{ij}\;\coloneqq\;\bm{T}_m\bD_{ij}\in\R^{D_m\times D},
\end{equation}
and its rescaled transpose used for the inverse rotation is
\begin{equation}\label{eq:wignerD-trunc-inv}
\widetilde\bD_{ij}^{-1}\;\coloneqq\;\widetilde\bD_{ij}^{\top}\,\circ\,\bm{S},\qquad
[\bm{S}]_{l,l'}\;=\;\sqrt{(2l+1)/(2\min(l,\mmax)+1)}\quad\text{for }l>\mmax,
\end{equation}
where $\circ$ denotes entrywise scaling by $\bm{S}$, and the rescaling restores the unit-norm normalisation lost by truncating high-$|m|$ components (see \texttt{so3.py:pre\_compute\_rotate\_inv\_rescale}).

\paragraph{Rotation-mask stabilisation.}
When $\bxi$ is within $10^{-6}$ of $\pm\by$, the Euler decomposition is singular.
For the gradient-prediction path the code detaches $\alpha,\gamma,\beta$ on those edges from the backward graph (\texttt{use\_rotation\_mask=True}); for the direct-prediction path the rotation is also detached from the graph entirely after construction.
This is purely numerical; it does not affect the symmetry properties of the forward pass.

% =============================================================
\section{Radial Basis and Smooth Cutoff}\label{sec:rad}

\subsection{Gaussian smearing}\label{sec:gauss}

The scalar edge distance $r_{ij}$ is expanded in $K_r=600$ Gaussian basis functions evenly spaced on $[0,\rcut]$:
\begin{equation}\label{eq:gauss}
\bigl[\eta(r)\bigr]_k\;=\;\exp\!\Bigl(-\tfrac{1}{2\sigma_g^2}(r-\mu_k)^2\Bigr),\quad
\mu_k=\rcut\,\frac{k-1}{K_r-1},\quad k=1,\dots,K_r,
\end{equation}
with $\sigma_g=2\,\Delta\mu$ where $\Delta\mu=\rcut/(K_r-1)$.
The output of \eqref{eq:gauss} is concatenated, on each edge, with the source-species and target-species embeddings $a^{\mathrm{src}}(Z_i),a^{\mathrm{tgt}}(Z_j)\in\R^{C_e}$ (with $C_e=128$ the edge-channel width).
The result $\bxi^{(e)}\in\R^{K_r+2C_e}$ is the edge-invariant feature consumed by every edge-conditioned scalar function in the architecture.

\subsection{Radial function block}\label{sec:radfn}

A \emph{radial function} (class \texttt{RadialFunction}) is a small MLP
\begin{equation}\label{eq:radfn}
\Hb_{\mathrm{rad}}\colon\R^{K_r+2C_e}\;\to\;\R^{C_{\mathrm{out}}},\qquad
\Hb_{\mathrm{rad}}(\bxi)\;=\;\mathrm{Linear}_2\!\bigl(\SiLU(\LN(\mathrm{Linear}_1(\bxi)))\bigr),
\end{equation}
with two hidden linear layers, layer normalisation, and $\SiLU$ activation in between.
Each callsite instantiates its own copy of $\Hb_{\mathrm{rad}}$ with independent learnable weights and a callsite-specific output width $C_{\mathrm{out}}$.
We denote callsite instances by a superscript:
\begin{equation}\label{eq:radfn-instances}
\Hb_{\mathrm{rad}}^{\mathrm{deg}}\;\text{(edge-degree embedding, \S\ref{sec:edgedeg})},\qquad
\Hb_{\mathrm{rad}}^{\mathrm{attn}}\;\text{(attention sublayer, \S\ref{sec:attn})}.
\end{equation}
A boolean flag \texttt{use\_rad\_l\_parametrization} controls the expansion mode of each instance: when \texttt{True} (the default) the per-edge radial vector is laid out per-$l$, identifying all $m$ within a fixed $l$; when \texttt{False} it is laid out per-$(l,|m|)$ stratum.

\subsection{Polynomial envelope}\label{sec:env}

A smooth cutoff is enforced by a polynomial envelope (class \texttt{PolynomialEnvelope})
\begin{equation}\label{eq:env}
\varphi(r)\;=\;
\begin{cases}
1\;+\;a\,(r/\rcut)^p\;+\;b\,(r/\rcut)^{p+1}\;+\;c\,(r/\rcut)^{p+2}, & r<\rcut,\\
0,&r\ge\rcut,
\end{cases}
\end{equation}
with $p=5$ and
\begin{equation}\label{eq:env-coefs}
a=-\tfrac{(p+1)(p+2)}{2},\qquad b=p(p+2),\qquad c=-\tfrac{p(p+1)}{2}.
\end{equation}
These coefficients are chosen so that $\varphi(0)=1,\varphi(\rcut)=0$, $\varphi'(\rcut)=\varphi''(\rcut)=0$.
The envelope is applied multiplicatively to the attention logits (\S\ref{sec:attn}) and to the edge-degree radial features (\S\ref{sec:edgedeg}).

% =============================================================
\section{Feature Spaces and Initial Embedding}\label{sec:init}

\paragraph{Feature space.}
The per-atom feature $\bh_i\in\R^{D\times C}$ is indexed by the spherical-harmonic linear address $\iota(l,m)=l^2+l+m$ on the first axis and by a channel index $c\in\{1,\dots,C\}$ on the second.
Under a global rotation $g\in\SO(3)$ the feature transforms as
\begin{equation}\label{eq:feature-action}
\bigl(g\cdot\bh_i\bigr)_{\iota(l,m),c}\;=\;\sum_{m'}\bigl[D^{(l)}(g)\bigr]_{m,m'}\,\bigl(\bh_i\bigr)_{\iota(l,m'),c},
\end{equation}
i.e.\ all channels share the rotational structure but transform independently of one another.

\paragraph{Species-only initial value.}
The node feature is initialised with the atom-type embedding $a(Z_i)\in\R^C$ in the $l=0$ slot and zero elsewhere:
\begin{equation}\label{eq:init-h}
\bh^{(0)}_i\;\coloneqq\;\bigl[a(Z_i),\;0,\;\dots,\;0\bigr]\in\R^{D\times C},
\end{equation}
where $a\colon\{1,\dots,Z_{\max}\}\to\R^C$ is a learned embedding table (\texttt{sphere\_embedding}).
The transformation rule \eqref{eq:feature-action} is consistent with this initialisation because $l=0$ features are invariant under rotation.

\subsection{Edge-degree embedding}\label{sec:edgedeg}

To populate the $l>0$ slots of $\bh^{(0)}_i$ before any transformer block runs, EquiformerV3 performs a single one-shot equivariant aggregation:
\begin{enumerate}[leftmargin=2em]
\item For each edge $(i,j)\in\Nb$, evaluate the edge-invariant feature $\bxi^{(e)}_{ij}$ of \S\ref{sec:gauss} and feed it through the edge-degree radial function $\Hb_{\mathrm{rad}}^{\mathrm{deg}}$ of \eqref{eq:radfn-instances} with output width $C\cdot(\lmax+1)$.
The output $\Hb_{\mathrm{rad}}^{\mathrm{deg}}(\bxi^{(e)}_{ij})$ is reshaped to $\bm{\rho}^{\mathrm{deg}}_{ij}\in\R^{(\lmax+1)\times C}$; this is a per-$l$, per-channel scalar.
\item Multiply by the envelope $\varphi(r_{ij})$.
\item Lift from $m=0$ to all $(l,m)$ by left-multiplying with the first $(\lmax+1)$ columns of the inverse-rotation Wigner $\widetilde\bD_{ij}^{-1}$ of \eqref{eq:wignerD-trunc-inv}; equivalently,
\begin{equation}\label{eq:edgedeg-lift}
\bm{e}^{\mathrm{deg}}_{ij}\;\coloneqq\;\bigl[\widetilde\bD_{ij}^{-1}\bigr]_{:,1:\lmax+1}\,\bm{\rho}^{\mathrm{deg}}_{ij}\;\in\;\R^{D\times C}.
\end{equation}
By construction, $\bm{e}^{\mathrm{deg}}_{ij}$ is the $\SO(3)$-equivariant lift of the radial $m=0$ stream by the inverse of the edge frame; rotating the edge rotates $\bm{e}^{\mathrm{deg}}_{ij}$.
\item Aggregate at the target with degree-mean normalisation:
\begin{equation}\label{eq:edgedeg-agg}
\Delta\bh^{\mathrm{deg}}_i\;\coloneqq\;\frac{1}{\overline d}\sum_{(i,j)\in\Nb}\bm{e}^{\mathrm{deg}}_{ij},
\end{equation}
where $\overline d$ is a fixed dataset-derived average degree (\texttt{avg\_degree}, equal to $23.4$ in the public IS2RE checkpoint).
\item Update the node feature:
\begin{equation}\label{eq:edgedeg-add}
\bh_i\;\leftarrow\;\bh^{(0)}_i\;+\;\Delta\bh^{\mathrm{deg}}_i.
\end{equation}
\end{enumerate}
After this step, $\bh_i\in\R^{D\times C}$ has non-zero content in every $(l,m)$ slot whenever the local geometry breaks isotropy, providing all subsequent transformer blocks with a non-trivial seed in the higher-$l$ slots.

% =============================================================
\section{\texorpdfstring{$\SO(2)_{\by}$ Linear Operators}{SO(2)\_y Linear Operators}}\label{sec:so2lin}

In the truncated $m$-primary layout (\S\ref{sec:mprim}) the action of $\SO(2)_{\by}\subset\SO(3)$ is block-diagonal in $|m|$.
EquiformerV3 builds two flavours of linear operator that respect this block structure.

\subsection{\texorpdfstring{The $m$-stratum linear operator $L_m$}{The m-stratum linear operator L\_m}}\label{sec:Lm}

For each $m\in\{1,\dots,\mmax\}$, the $m$-stratum carries $2\bigl(\lmax-m+1\bigr)$ real coordinates: the real and imaginary parts of the $(l,\pm m)$ pairs for $l=m,m+1,\dots,\lmax$.
$\SO(2)_{\by}$ acts on $(x^+_m,x^-_m)\in\R^2$ by rotation, equivalently by complex multiplication on $x^+_m+i\,x^-_m$.
Schur for $\SO(2)_{\by}$ over $\R$ implies that the $\SO(2)_{\by}$-equivariant $\R$-linear maps $\R^{2(\lmax-m+1)}\otimes\R^{C_{\mathrm{in}}}\;\to\;\R^{2(\lmax-m+1)}\otimes\R^{C_{\mathrm{out}}}$ are parameterised by a single \emph{complex} matrix
\begin{equation}\label{eq:Lm-cmat}
\bW_m\;=\;\bU_m\;+\;i\,\bV_m\;\in\;\C^{(\lmax-m+1)C_{\mathrm{out}}\times(\lmax-m+1)C_{\mathrm{in}}}.
\end{equation}
Implementationally, $\bU_m,\bV_m$ are obtained from a single real linear layer with $2(\lmax-m+1)C_{\mathrm{out}}$ outputs followed by a split.
The operator $L_m$ acts on $(x^+,x^-)$ by the real form of complex multiplication:
\begin{equation}\label{eq:Lm-action}
\begin{pmatrix} L_m x \\[2pt] L_m x\end{pmatrix}^+ \;=\; \bU_m\,x^+\;-\;\bV_m\,x^-,\qquad
\begin{pmatrix} L_m x \\[2pt] L_m x\end{pmatrix}^- \;=\; \bU_m\,x^-\;+\;\bV_m\,x^+.
\end{equation}
The weight is uniformly initialised and then rescaled by $1/\sqrt{2}$ so that the output variance matches that of a real layer of the same effective fan-in.

\subsection{\texorpdfstring{The $m=0$ scalar linear $L_0$}{The m=0 scalar linear L\_0}}\label{sec:L0}

The $m=0$ stratum carries one real coordinate per $l$, i.e.\ $(\lmax+1)C_{\mathrm{in}}$ real degrees of freedom.
$L_0$ is an unconstrained dense linear map with a bias term:
\begin{equation}\label{eq:L0}
L_0\colon\R^{(\lmax+1)C_{\mathrm{in}}}\;\to\;\R^{(\lmax+1)C_{\mathrm{out}}+C^{\mathrm{ex}}},\qquad
y\;=\;\bA_0\,x+\bm{b}_0,
\end{equation}
with optional \emph{extra} scalar outputs of width $C^{\mathrm{ex}}$ used as side-channels for attention logits and grid-activation gates (\S\ref{sec:attn}, \S\ref{sec:ffn}).

\subsection{\texorpdfstring{Compound $\SO(2)_{\by}$ linear operator $L^{\mathrm{SO2}}$}{Compound SO(2)\_y linear operator}}\label{sec:Lso2}

The full $\SO(2)_{\by}$-equivariant linear operator on a truncated feature is
\begin{equation}\label{eq:LSO2}
L^{\mathrm{SO2}}(x)\;\coloneqq\;\bigl(\,L_0(x_0)\,,\;L_1(x_1)\,,\;\dots\,,\;L_{\mmax}(x_{\mmax})\,\bigr),
\end{equation}
acting independently on each $|m|$-stratum.
The shape conventions are: input $x\in\R^{D_m\times C_{\mathrm{in}}}$, output $\bigl(\,\R^{(\lmax+1)\times C_{\mathrm{out}}+C^{\mathrm{ex}}},\;\R^{2(\lmax-1+1)\times C_{\mathrm{out}}},\;\dots,\;\R^{2(\lmax-\mmax+1)\times C_{\mathrm{out}}}\bigr)$.
The extra-scalar block of $L_0$ is split out and returned separately by the implementation (\texttt{so2\_ops.py:SO2Linear}).

\begin{proposition}[$\SO(2)_{\by}$ equivariance of $L^{\mathrm{SO2}}$]\label{prop:Lso2}
For any $g_\theta=R_y(\theta)\in\SO(2)_{\by}$ acting on truncated features in the $m$-primary layout by the block-diagonal matrix $\bD^m(\theta)=\diag\bigl(I,\,R^+_1(\theta),\dots,R^+_{\mmax}(\theta)\bigr)$, we have
\begin{equation}\label{eq:LSO2-equiv}
L^{\mathrm{SO2}}\bigl(\bD^m(\theta)\,x\bigr)\;=\;\bD^m(\theta)\,L^{\mathrm{SO2}}(x).
\end{equation}
This follows from \eqref{eq:Lm-action} and the fact that $L_0$ is the identity in the $m$-action.
\end{proposition}

\begin{remark}[Equivalence with DPA4's $\SO(2)$-linear operator]\label{rmk:so2-equiv}
The operator $L^{\mathrm{SO2}}$ of \eqref{eq:LSO2} is mathematically the same operator as DPA4's $\SO(2)$-linear mixer (archi\_dpa4.tex, \S\ref{sec:ops}; source: \texttt{sezm\_nn/so2.py:SO2Linear}), up to the choice of fixed axis ($\by$ here vs.\ $\bz$ in DPA4) and the absence of the multi-stream \emph{focus} axis.
By Schur's lemma applied to the real $\SO(2)$ action on the $m$-major reduced layout, the space of $\R$-linear $\SO(2)$-equivariant maps decomposes as
\begin{equation}\label{eq:hom-so2-eqv3}
\mathrm{Hom}_{\SO(2)}\bigl(\R^{D_m}\otimes\R^{C_{\mathrm{in}}},\;\R^{D_m}\otimes\R^{C_{\mathrm{out}}}\bigr)
\;=\;\R^{(\lmax+1)C_{\mathrm{out}}\times(\lmax+1)C_{\mathrm{in}}}\;\bigoplus_{m=1}^{\mmax}\C^{(\lmax-m+1)C_{\mathrm{out}}\times(\lmax-m+1)C_{\mathrm{in}}},
\end{equation}
where the first summand is the unconstrained real matrix on the $m=0$ stratum and each complex summand encodes the $2\times 2$ real block
\begin{equation}\label{eq:so2-block}
\begin{pmatrix}\bU_m & -\bV_m\\ \bV_m & \bU_m\end{pmatrix},\qquad \bW_m\;=\;\bU_m+i\bV_m\in\C^{(\lmax-m+1)C_{\mathrm{out}}\times(\lmax-m+1)C_{\mathrm{in}}}.
\end{equation}
Both architectures parameterise the full space \eqref{eq:hom-so2-eqv3}: cross-$l$ mixing within a fixed $|m|$-stratum is permitted (the blocks $\bU_m,\bV_m$ are dense in their $(\lmax-m+1)$ row/column directions), cross-$|m|$ mixing is forbidden (block-diagonality in $|m|$).
Both apply the variance-preserving rescale $1/\sqrt{2}$ to the $|m|>0$ weights to compensate for the doubled real parameter count of the $2\times 2$ block.
The implementations differ in incidentals only:
\begin{itemize}[leftmargin=2em]
\item \emph{Multi-stream.}
DPA4 carries $F$ independent in-frame focus streams with independent $\bU_m,\bV_m$ per focus, assembled into a single block-diagonal $\R^{D_m C_{\mathrm{in}}\times F\times D_m C_{\mathrm{out}}}$ weight and applied via one batched matmul (\texttt{einsum("efi,ifo->efo")}); EquiformerV3 uses $F=1$ and loops over $|m|$ to concatenate the per-stratum outputs.
\item \emph{Side-channel scalars.}
EquiformerV3 attaches additional scalar outputs to $L_0$ via the \texttt{extra\_m0\_out\_channels} argument; these are $\SO(2)$-invariant by construction and feed attention logits and grid-activation gates (\S\ref{sec:attn-soft}, \eqref{eq:attn-act}).
DPA4 keeps the analogous scalar side-channels in separate modules (\texttt{attn\_res.py}, \texttt{attention.py}) rather than as extra $L_0$ outputs.
The choice has no effect on the equivariance class of $L^{\mathrm{SO2}}$.
\item \emph{Bias.}
EquiformerV3's $L_0$ has a plain bias on the $l=0$ output.
DPA4's $L_0$ has an optional bias coupled to the radial--envelope product through the correction $b_0\bigl(\widetilde\rho_{l=0}\,\varphi(r)-1\bigr)$ so that the bias contribution vanishes whenever the edge contribution vanishes (relevant for the source-freeze bridging zone).
This is a feature of the surrounding edge-convolution wrapper, not of the $\SO(2)$-equivariant linear operator class.
\item \emph{Weight assembly and caching.}
DPA4 assembles the full block-diagonal weight once per forward pass and caches it in evaluation mode (\texttt{so2.py:\_build\_so2\_weight}); EquiformerV3 keeps the strata in separate \texttt{nn.Linear} layers and applies them sequentially.
\item \emph{Initialisation.}
DPA4 uses truncated-normal fan-in/out; EquiformerV3 uses uniform $[-1/\sqrt{n},+1/\sqrt{n}]$.
Both belong to the Glorot family and produce the same expected output variance to leading order.
\end{itemize}
None of these incidentals changes the equivariant content: the two architectures realise the same $\mathrm{Hom}_{\SO(2)}$ space \eqref{eq:hom-so2-eqv3} and any operator producible in one is producible in the other (with $F=1$ on the DPA4 side).
What does differ between the architectures, and substantively, is how this operator is \emph{used}: how many times it is stacked, what nonlinearity is sandwiched between stacks, whether the aggregation downstream is sum-over-edges or softmax attention, and whether the per-edge message is built from $\bh_j$ alone or from $[\bh_i\|\bh_j]$.
These higher-level choices are the subject of \S\ref{sec:compare-diff}.
\end{remark}

% =============================================================
\section{$\Sph$-Grid Activation and the FFN}\label{sec:ffn}

The FFN needs a nonlinearity that respects $\SO(3)$-equivariance.
EquiformerV3 implements one by a three-step sandwich: \emph{sample} the irrep feature at a fixed set of points on $\Sph$, apply an unconstrained pointwise MLP at every sample point, and \emph{project} the resulting samples back to irrep coefficients.
The purpose of this section is to construct the sandwich and to give a complete proof of its equivariance.
The non-trivial fact is that the grid values are \emph{not} $\SO(3)$-invariant: they transform under rotation by a row-permutation of grid points.
Equivariance is recovered because the pointwise MLP commutes with that row-permutation, and the pull-back converts the row-permutation back into the Wigner-$D$ action on irrep coefficients.

\subsection{The $\Sph$-grid}\label{sec:S2grid}

Fix a latitude resolution $N_\beta$ and a longitude resolution $N_\alpha$ (defaults: $N_\beta=20$, $N_\alpha=8$ for the attention sublayer's activation, $N_\beta=N_\alpha=20$ for the FFN's activation).
The grid is the tensor product
\begin{equation}\label{eq:grid}
\mathcal{G}\;=\;\bigl\{\hat\br_{ab}\in\Sph\,:\,a=1,\dots,N_\beta,\;b=1,\dots,N_\alpha\bigr\},
\end{equation}
where $\hat\br_{ab}$ is the unit vector with spherical coordinates $(\beta_a,\alpha_b)$, $\beta_a$ are Gauss--Legendre nodes on $[0,\pi]$, and $\alpha_b=2\pi(b-1)/N_\alpha$ are equispaced on $[0,2\pi)$.
The grid is \emph{fixed in the lab frame} for the duration of the architecture; it is not attached to any atom or edge.

\subsection{Sample: irrep coefficients become grid values}\label{sec:topull}

Define the sample matrix $\bm{S}\in\R^{(N_\beta N_\alpha)\times D}$ row by row:
\begin{equation}\label{eq:sample}
[\bm{S}]_{(a,b),(l,m)}\;=\;s_l\,Y_l^m(\hat\br_{ab}),
\end{equation}
where $Y_l^m$ is the real spherical harmonic and $s_l$ is a rescale factor restoring unit-norm normalisation when the truncated $\mmax$ layout is used ($s_l=\sqrt{(2l+1)/(2\min(l,\mmax)+1)}$ for $l>\mmax$, and $s_l=1$ otherwise).
For an irrep feature $\bh\in\R^{D\times C}$ define the grid values
\begin{equation}\label{eq:grid-eval}
\bh^{\Sph}\;\coloneqq\;\bm{S}\,\bh\;\in\;\R^{(N_\beta N_\alpha)\times C}.
\end{equation}
The row $(a,b)$ of $\bh^{\Sph}$ is the value at $\hat\br_{ab}$ of the channel-wise band-limited function
\begin{equation}\label{eq:fh}
f^{c}_{\bh}(\hat\br)\;\coloneqq\;\sum_{l\le\lmax,\,m}\bh_{(l,m),c}\,Y_l^m(\hat\br),\qquad
[\bh^{\Sph}]_{(a,b),c}\;=\;f^{c}_{\bh}(\hat\br_{ab}).
\end{equation}
The map $\bh\mapsto f_{\bh}$ is a linear isomorphism between the irrep space $V_{\le\lmax}$ and the space of band-limited functions on $\Sph$; the matrix $\bm{S}$ is the discretisation of this isomorphism on the grid.

\subsection{Pull-back: grid values become irrep coefficients}\label{sec:pull}

Define the pull-back $\bm{S}^{\dagger}\in\R^{D\times(N_\beta N_\alpha)}$ as the discrete spherical-harmonic projection: applied to a function on the grid, it returns the irrep coefficients of the band-limited reconstruction.
The construction in \texttt{e3nn.FromS2Grid} uses the same Gauss--Legendre quadrature so that
\begin{equation}\label{eq:Sdag-S}
\bm{S}^{\dagger}\,\bm{S}\;=\;I_D\qquad\text{whenever the quadrature is exact at the relevant band limit.}
\end{equation}
At the FFN's default $(N_\beta,N_\alpha)=(20,20)$ and $\lmax=6$, the quadrature is exact for spherical-harmonic products up to the band limit produced by the pointwise MLP on $\lmax$-band-limited inputs, so \eqref{eq:Sdag-S} holds and the pull-back is a left inverse of the sample.

\subsection{How rotation acts on grid values}\label{sec:rot-on-grid}

This is the step on which the equivariance argument turns.
Under a global rotation $g\in\SO(3)$, the configuration rotates, and with it the irrep feature transforms by Wigner-$D$: $\bh\mapsto\bD(g)\bh$.
The grid points $\hat\br_{ab}$ stay fixed in the lab frame.
Tracking through the sample using the spherical-harmonic transformation law $Y_l^m(g\hat\br)=\sum_{m'}D^{(l)}_{m,m'}(g)Y_l^{m'}(\hat\br)$ (equivalently, $\bD(g)^{\top}\bm{Y}_{ab}=\bm{Y}_{g^{-1}\hat\br_{ab}}$ where $\bm{Y}_{ab}$ is the row vector of $\bm{S}$ at index $(a,b)$):
\begin{equation}\label{eq:rot-grid}
[\bm{S}\,(\bD(g)\bh)]_{(a,b),c}\;=\;\langle\bm{Y}_{ab},\bD(g)\bh_{:,c}\rangle\;=\;\langle\bD(g)^{\top}\bm{Y}_{ab},\bh_{:,c}\rangle\;=\;\langle\bm{Y}_{g^{-1}\hat\br_{ab}},\bh_{:,c}\rangle\;=\;f^{c}_{\bh}(g^{-1}\hat\br_{ab}).
\end{equation}
\emph{In words}: after rotating the feature, the grid value at index $(a,b)$ is the value of the \emph{original} function at the \emph{back-rotated} grid point $g^{-1}\hat\br_{ab}$.
The grid values are not invariant — they are permuted (or, for a non-aligned $g$, interpolated) across grid indices.

When the grid is closed under $g^{-1}$, i.e.\ when there exists a permutation $\pi_g$ of grid indices such that $g^{-1}\hat\br_{ab}=\hat\br_{\pi_g(a,b)}$, equation \eqref{eq:rot-grid} reads
\begin{equation}\label{eq:rot-perm}
\bm{S}\,(\bD(g)\bh)\;=\;P_g\,(\bm{S}\bh),
\end{equation}
where $P_g$ is the $(N_\beta N_\alpha)\times(N_\beta N_\alpha)$ row-permutation matrix with $[P_g]_{(a,b),(a',b')}=\delta_{(a',b'),\pi_g(a,b)}$.
For a generic $g$ the closure fails and $P_g$ is replaced by a band-limited interpolation operator; the continuous statement \eqref{eq:rot-grid} is exact in both cases, only the discrete representation \eqref{eq:rot-perm} requires closure.

The takeaway from \eqref{eq:rot-grid}--\eqref{eq:rot-perm}: \emph{$\bm{S}$ intertwines the Wigner-$D$ action on irrep coefficients with the function-shift action on grid samples.}
Symbolically, $\bm{S}\circ\bD(g)\;=\;P_g\circ\bm{S}$.

\subsection{Pull-back intertwines back}\label{sec:pull-intertwines}

The reverse statement holds for $\bm{S}^{\dagger}$.
For a function on the grid, the rotation action is row-permutation $P_g$ (or interpolation); the spherical-harmonic projection of the rotated function is the Wigner-$D$-rotation of the projection.
Concretely, applying \eqref{eq:Sdag-S} and \eqref{eq:rot-perm},
\begin{equation}\label{eq:Sdag-intertwine}
\bm{S}^{\dagger}\,P_g\;=\;\bD(g)\,\bm{S}^{\dagger},
\end{equation}
i.e.\ $\bm{S}^{\dagger}$ intertwines the function-shift action back into the Wigner-$D$ action.
This is just the adjoint of the statement $\bm{S}\circ\bD(g)=P_g\circ\bm{S}$ from the previous subsection, using that $\bm{S}^{\dagger}\bm{S}=I_D$ on the band-limited subspace.

\subsection{Pointwise channel-mixing commutes with row-permutation}\label{sec:point-comm}

Call a map $\Phi\colon\R^{(N_\beta N_\alpha)\times C}\to\R^{(N_\beta N_\alpha)\times C}$ \emph{pointwise in the grid index} if there is a single function $\phi\colon\R^C\to\R^C$ such that, for every grid index $(a,b)$,
\begin{equation}\label{eq:Phi-pointwise}
[\Phi(X)]_{(a,b),:}\;=\;\phi(X_{(a,b),:}).
\end{equation}
In words: $\Phi$ does not see the grid index; it sees only the channel vector at each row, and the \emph{same} channel-mixing rule applies at every row.
The grid MLP \eqref{eq:gridmlp} below is exactly such a $\Phi$.

\begin{lemma}[Pointwise commutes with permutation]\label{lem:point-perm}
For any permutation matrix $P$ on grid indices and any $\Phi$ of the form \eqref{eq:Phi-pointwise},
\begin{equation}\label{eq:point-perm}
\Phi(P\,X)\;=\;P\,\Phi(X).
\end{equation}
\end{lemma}

\begin{proof}
Both sides have row $(a,b)$ equal to $\phi(X_{P(a,b),:})$.
The left side gives this by definition of $\Phi$ applied to the permuted matrix; the right side gives this because $\Phi$ acts row-by-row with the same $\phi$, and the permutation acts last on rows.
\end{proof}

This is the entire mechanism by which a pointwise nonlinearity preserves equivariance.
The crucial structural property is that $\phi$ does not depend on the grid index — it is the same function at every row.
Therefore moving rows around (i.e.\ rotating) before or after $\Phi$ gives the same result.

\subsection{The equivariance theorem}\label{sec:eqv-thm}

\begin{theorem}[Equivariance of the grid sandwich]\label{thm:grid-sandwich}
Let $\Phi\colon\R^{(N_\beta N_\alpha)\times C}\to\R^{(N_\beta N_\alpha)\times C}$ be any map pointwise in the grid index, as in \eqref{eq:Phi-pointwise}.
Suppose the quadrature is exact at the band limit produced by $\Phi$ on $\lmax$-band-limited inputs, so \eqref{eq:Sdag-S} holds.
Then the grid sandwich
\begin{equation}\label{eq:sandwich}
\mathcal{F}\;\coloneqq\;\bm{S}^{\dagger}\,\circ\,\Phi\,\circ\,\bm{S}\;\colon\;\R^{D\times C}\longrightarrow\R^{D\times C}
\end{equation}
is $\SO(3)$-equivariant: for every $g\in\SO(3)$ and every $\bh\in\R^{D\times C}$,
\begin{equation}\label{eq:sandwich-equiv}
\mathcal{F}(\bD(g)\bh)\;=\;\bD(g)\,\mathcal{F}(\bh).
\end{equation}
\end{theorem}

\begin{proof}
Apply $\mathcal{F}$ to the rotated feature and use the three intertwining steps developed above:
\begin{align*}
\mathcal{F}(\bD(g)\bh)
\;&=\;\bm{S}^{\dagger}\,\Phi\,\bm{S}(\bD(g)\bh)\\
\;&\stackrel{\eqref{eq:rot-perm}}{=}\;\bm{S}^{\dagger}\,\Phi\,(P_g\,\bm{S}\bh)\\
\;&\stackrel{\eqref{eq:point-perm}}{=}\;\bm{S}^{\dagger}\,P_g\,\Phi(\bm{S}\bh)\\
\;&\stackrel{\eqref{eq:Sdag-intertwine}}{=}\;\bD(g)\,\bm{S}^{\dagger}\,\Phi(\bm{S}\bh)\\
\;&=\;\bD(g)\,\mathcal{F}(\bh).
\end{align*}
The first equality is the definition of $\mathcal{F}$.
The second uses that $\bm{S}$ intertwines Wigner-$D$ on irreps with permutation on the grid (\S\ref{sec:rot-on-grid}).
The third uses that $\Phi$ pointwise commutes with row-permutation (Lemma~\ref{lem:point-perm}).
The fourth uses that $\bm{S}^{\dagger}$ intertwines permutation on the grid back into Wigner-$D$ on irreps (\S\ref{sec:pull-intertwines}).
For generic $g$ that does not close the grid, replace $P_g$ by the band-limited interpolation operator and replace \eqref{eq:rot-perm}--\eqref{eq:Sdag-intertwine} by their function-space versions $\bm{S}\circ\bD(g)=R_g\circ\bm{S}$ and $\bm{S}^{\dagger}\circ R_g=\bD(g)\circ\bm{S}^{\dagger}$, where $R_g$ is the function-shift $f\mapsto f\circ g^{-1}$; the same three-line chain goes through verbatim.
\end{proof}

The theorem makes precise what was loosely stated earlier: equivariance does \emph{not} follow from the grid values being invariant — they are not.
It follows from the conjunction of three facts: the sample intertwines Wigner-$D$ on irreps with the rotation action on grid samples; pointwise channel-mixing commutes with that rotation action; the pull-back intertwines back.

\subsection{Two non-negotiable design constraints}\label{sec:constraints}

Theorem~\ref{thm:grid-sandwich} relies on $\Phi$ being pointwise in the grid index.
Two architectural choices are therefore essential:

\begin{enumerate}[leftmargin=2em]
\item \emph{No mixing across grid points.}
Any weight that couples grid index $(a_1,b_1)$ to $(a_2,b_2)$ — a $1\times 1$ convolution across $(a,b)$, an attention layer over grid points, a pooling over the grid — would violate the pointwise condition \eqref{eq:Phi-pointwise} and break the commutation with $P_g$.
The grid MLP is therefore implemented as a per-row channel mixer with the same weights at every row.
\item \emph{Any scalar gate is a $\SO(3)$-invariant of $\bh$.}
A gate that depends on the grid index $(a,b)$ would assign different multipliers to different grid points, violating the pointwise condition.
A gate that depends on a non-$l=0$ slot of $\bh$ would itself rotate under $g$, again violating the condition.
The grid MLP restricts the gate to depend only on $\bh_{l=0}$, which is rotation-invariant, so the same scalar gate multiplies every grid point in every rotated copy of the configuration.
\end{enumerate}

Both constraints are satisfied by the grid MLP defined below.

\subsection{Grid MLP}\label{sec:gridmlp}

The grid MLP applied in the FFN has the form
\begin{equation}\label{eq:gridmlp}
\mathcal{M}_{\Sph}(\bh^{\Sph};\bm{s})\;=\;\bm{W}_2\bigl(\,\SiLU(\bm{W}_1\bh^{\Sph})\;\odot\;\Sigmoid(\bm{w}_g\bm{s})\,\bigr),
\end{equation}
with $\bm{W}_1\in\R^{C_h\times C}$, $\bm{W}_2\in\R^{C\times C_h}$, $\bm{w}_g\in\R^{C_h\times C}$ learnable, $\bm{s}=\bh_{l=0}\in\R^C$ the rotation-invariant scalar side-channel, and $\odot$ pointwise multiplication on channels.
At each grid row $(a,b)$ the operation is the same channel-mixing recipe; no operator looks at the grid index.
In the default activation \texttt{sep-merge\_gates2\_swiglu} the $\SiLU$ in \eqref{eq:gridmlp} is replaced by SwiGLU,
\begin{equation}\label{eq:swiglu}
\SwiGLU(x;\bm{W}_a,\bm{W}_b)\;=\;(\bm{W}_a x)\odot\SiLU(\bm{W}_b x).
\end{equation}
By Theorem~\ref{thm:grid-sandwich}, the inner sandwich $\bm{S}^{\dagger}\circ\mathcal{M}_{\Sph}\circ\bm{S}$ is $\SO(3)$-equivariant.

\subsection{Quadrature exactness and aliasing}\label{sec:quad}

The hypothesis \eqref{eq:Sdag-S} of Theorem~\ref{thm:grid-sandwich} requires the quadrature to be exact at the band limit produced by $\Phi$ on $\lmax$-band-limited inputs.
The Gauss--Legendre--Fourier quadrature on the e3nn grid is exact for the integral of any spherical-harmonic product $Y_{l_1}^{m_1}Y_{l_2}^{m_2}$ with $l_1+l_2\le 2L_{\mathrm{quad}}$, where $L_{\mathrm{quad}}$ is determined by $(N_\beta,N_\alpha)$.
For a SwiGLU activation acting on $\lmax$-band-limited inputs, the relevant band limit is $2\lmax$, and the FFN default $(N_\beta,N_\alpha)=(20,20)$ at $\lmax=6$ comfortably exceeds this.
If the grid were under-resolved, the pointwise nonlinearity $\phi$ would generate angular frequencies above the quadrature band-limit; the discrete pull-back $\bm{S}^{\dagger}$ would alias them onto the kept $l\le\lmax$ coefficients in a $g$-dependent way, and the third equality in the proof of Theorem~\ref{thm:grid-sandwich} would hold only up to a discretisation error.

\subsection{The equivariant FFN}\label{sec:effn}

The full FFN sublayer is the channel-mixed grid sandwich
\begin{equation}\label{eq:effn-1}
\bh\;\xrightarrow{L_1^{\mathrm{ch}}}\;\widetilde\bh\;\xrightarrow{\bm{S}}\;\widetilde\bh^{\Sph}\;\xrightarrow{\mathcal{M}_{\Sph}(\cdot;\bm{s})}\;\bg\;\xrightarrow{\bm{S}^{\dagger}}\;\bg^{\mathrm{irr}}\;\xrightarrow{L_2^{\mathrm{ch}}}\;\Delta\bh^{\mathrm{ffn}},
\end{equation}
with $\bm{s}=\bh_{l=0}$.
Here $L_1^{\mathrm{ch}},L_2^{\mathrm{ch}}$ are \emph{degree-preserving channel mixers} implemented by \texttt{SO3Linear}: for each $l$ there is one $C_h\times C$ weight matrix shared across all $2l+1$ values of $m$.
These are $\SO(3)$-equivariant because they act in the global frame, mix only the channel axis at fixed $l$, and treat all $m$ uniformly (so they commute with each $D^{(l)}(g)$).
Combined with Theorem~\ref{thm:grid-sandwich} for the inner sandwich, the entire FFN sublayer $\Hb^{\mathrm{ffn}}_\ell=L_2^{\mathrm{ch}}\circ\mathcal{F}\circ L_1^{\mathrm{ch}}$ commutes with $\bD(g)$, hence is $\SO(3)$-equivariant.

A \emph{merge} variant adds a separate scalar MLP output to the $l=0$ slot of $\bg^{\mathrm{irr}}$:
\begin{equation}\label{eq:scalar-merge}
\bigl[\bg^{\mathrm{irr}}\bigr]_{l=0}\;\leftarrow\;\bigl[\bg^{\mathrm{irr}}\bigr]_{l=0}\;+\;\mathrm{Linear}(\SwiGLU(\bm{s})).
\end{equation}
The added term depends only on $\bh_{l=0}=\bm{s}$ (an $\SO(3)$-invariant) and is written into the $l=0$ slot (also $\SO(3)$-invariant), so the merge is $\SO(3)$-equivariant.
No per-edge Wigner-$D$ rotation, $\SO(2)$-equivariant operator, or $\mmax$ truncation appears anywhere in the FFN: the FFN sees the full global-frame irrep feature $\bh\in\R^{D\times C}$ throughout.

% =============================================================
\section{Equivariant Graph Attention}\label{sec:attn}

The attention sublayer is where edge information enters: each edge produces a message in the local frame, the message is scaled by a softmax-attention weight, the messages are rotated back to the global frame, and the per-edge messages are summed at each target.

\subsection{Edge feature construction}\label{sec:attn-msg}

For each edge $(i,j)\in\Nb$, the attention sublayer first builds an edge feature in $\R^{D\times C_{\mathrm{m}}}$:
\begin{equation}\label{eq:attn-edge}
\bm{x}_{ij}^{\mathrm{cat}}\;\coloneqq\;\bigl[\,\bh_i^{(\text{src})}\;\big\|\;\bh_j^{(\text{tgt})}\,\bigr]\in\R^{D\times 2C},
\end{equation}
where $\|$ denotes channel concatenation.
(An alternative \texttt{use\_add\_merge=True} replaces the concatenation by a weighted sum; we describe the default concat branch and the variant in Remark~\ref{rmk:addmerge}.)
The radial scaling
\begin{equation}\label{eq:attn-rad}
\widetilde\brho_{ij}\;\coloneqq\;\Hb_{\mathrm{rad}}^{\mathrm{attn}}(\bxi^{(e)}_{ij})\in\R^{(\lmax+1)\times 2C}
\end{equation}
is broadcast across $m$ (one weight per $l$) and applied entrywise:
\begin{equation}\label{eq:attn-rad-apply}
\bm{x}_{ij}^{\mathrm{cat,rad}}\;=\;\widetilde\brho_{ij}\;\odot\;\bm{x}_{ij}^{\mathrm{cat}}.
\end{equation}
The radially-scaled feature is then rotated into the local frame:
\begin{equation}\label{eq:attn-rot}
\bm{x}_{ij}^{\mathrm{loc}}\;\coloneqq\;\widetilde\bD_{ij}\,\bm{x}_{ij}^{\mathrm{cat,rad}}\in\R^{D_m\times 2C}.
\end{equation}
(When the radial parametrisation flag is off, the rotation is performed first and the scaling, now per-$|m|$-stratum, is applied second; this is mathematically equivalent.)

\begin{remark}\label{rmk:addmerge}
The \texttt{use\_add\_merge} variant replaces \eqref{eq:attn-edge}--\eqref{eq:attn-rad-apply} by
$\bm{x}_{ij}^{\mathrm{loc}}=\widetilde\bD_{ij}(\widetilde\brho^{\mathrm{src}}_{ij}\odot\bh_i+\widetilde\brho^{\mathrm{tgt}}_{ij}\odot\bh_j)$, saving roughly half the Wigner multiplication cost; an explicit $1/\sqrt{2}$ rescaling on the radial-MLP head preserves variance.
\end{remark}

\subsection{In-frame projection and grid activation}\label{sec:attn-proj}

The local-frame feature is mixed by an $\SO(2)_{\by}$-equivariant linear operator that outputs both the value stream and the side-channel scalar features used for attention:
\begin{equation}\label{eq:so2-1}
\bigl(\bm{v}^{\mathrm{loc}}_{ij},\,\balpha_{ij}\bigr)\;\coloneqq\;L^{\mathrm{SO2}}_1\bigl(\bm{x}_{ij}^{\mathrm{loc}}\bigr),\qquad
\bm{v}^{\mathrm{loc}}_{ij}\in\R^{D_m\times C_h},\;\balpha_{ij}\in\R^{C^{\mathrm{ex}}}.
\end{equation}
The extra scalars $\balpha_{ij}$ are split into the attention-logit features $\bm{a}_{ij}\in\R^{H C^{\mathrm{a}}}$ and the gate scalars $\bm{s}_{ij}\in\R^{C_h+C_h/2}$ that drive the SwiGLU gates of the activation.
Because the extra-scalar block of $L_0$ acts on the $m=0$ stratum of $\bm{x}^{\mathrm{loc}}_{ij}$, which is $\SO(3)$-invariant (the global rotation is absorbed by the per-edge frame change) and trivially $\SO(2)_{\by}$-invariant (the $m=0$ irrep of $\SO(2)$ is the trivial one), both $\bm{a}_{ij}$ and $\bm{s}_{ij}$ are $\SO(3)$-invariant scalar tuples.
This invariance is what makes the attention logits below scalar-valued under rotation, and is what the equivariance of the whole attention sublayer (Proposition~\ref{prop:attn-equiv}) hinges on.
The value stream is passed through a separable $\Sph$-grid activation
\begin{equation}\label{eq:attn-act}
\bm{v}^{\mathrm{loc}}_{ij}\;\leftarrow\;\bm{S}^{\dagger}\,\bigl(\,\SwiGLU\bigl(\bm{S}\,\bm{v}^{\mathrm{loc}}_{ij}\bigr)\,\odot\,\Sigmoid(\bm{w}_g\bm{s}_{ij})\,\bigr),
\end{equation}
which is the same construction as \eqref{eq:gridmlp} but with a separate scalar path added back into the $l=0$ slot when the activation name carries the \texttt{merge} flag.
A second $\SO(2)_{\by}$-equivariant linear $L^{\mathrm{SO2}}_2$ then re-projects to the per-head value width:
\begin{equation}\label{eq:so2-2}
\bm{v}^{\mathrm{loc}}_{ij}\;\leftarrow\;L^{\mathrm{SO2}}_2\bigl(\bm{v}^{\mathrm{loc}}_{ij}\bigr)\in\R^{D_m\times H C^{\mathrm{v}}}.
\end{equation}

\subsection{Attention logits and softmax}\label{sec:attn-soft}

The attention logits per head are computed from $\bm{a}_{ij}\in\R^{H C^{\mathrm{a}}}$, reshaped to $(H,C^{\mathrm{a}})$:
\begin{equation}\label{eq:attn-logit}
\widetilde a^{(h)}_{ij}\;\coloneqq\;\bigl\langle\,\bm{q}^{(h)}\,,\;\mathrm{SmoothLeakyReLU}(\LN(\bm{a}^{(h)}_{ij}))\,\bigr\rangle,\qquad h=1,\dots,H,
\end{equation}
with a learnable query $\bm{q}^{(h)}\in\R^{C^{\mathrm{a}}}$ initialised uniformly at $\pm 1/\sqrt{C^{\mathrm{a}}}$, the layer norm acting on the $C^{\mathrm{a}}$-feature axis, and
$\mathrm{SmoothLeakyReLU}(x)=\bigl(x+\sqrt{x^2+\epsilon}\bigr)/2$ a smoothed leaky-ReLU with $\epsilon=10^{-4}$.
An optional bounded \emph{softcap} maps $\widetilde a^{(h)}_{ij}\mapsto\softcap\tanh(\widetilde a^{(h)}_{ij}/\softcap)$ to clip extreme values.

The softmax is applied per target with envelope rescaling:
\begin{equation}\label{eq:attn-softmax}
\alpha^{(h)}_{ij}\;\coloneqq\;\varphi(r_{ij})\,\frac{\varphi(r_{ij})\exp(\widetilde a^{(h)}_{ij})}{\eta_{\mathrm{eps}}+\sum_{(i',j)\in\Nb}\varphi(r_{i'j})\exp(\widetilde a^{(h)}_{i'j})},
\end{equation}
with $\eta_{\mathrm{eps}}=10^{-16}$ a numerical floor.
The double envelope multiplication is the architectural signature: it makes both the numerator and the denominator vanish smoothly as any edge reaches the cutoff, so that the per-edge attention weight is $C^2$-smooth at $r_{ij}=\rcut$.

\subsection{Aggregation and projection}\label{sec:attn-agg}

The attention-weighted local-frame value is computed per head and per channel:
\begin{equation}\label{eq:attn-weight}
\bm{u}^{\mathrm{loc}}_{ij}\;\coloneqq\;\alpha_{ij}\;\odot\;\bm{v}^{\mathrm{loc}}_{ij}\in\R^{D_m\times H C^{\mathrm{v}}},
\end{equation}
where the head index $h$ is broadcast over the channel block of width $C^{\mathrm{v}}$.
The local-frame message is then rotated back to the global frame
\begin{equation}\label{eq:attn-back}
\bm{u}_{ij}\;\coloneqq\;\widetilde\bD_{ij}^{-1}\,\bm{u}^{\mathrm{loc}}_{ij}\in\R^{D\times H C^{\mathrm{v}}},
\end{equation}
summed at the target,
\begin{equation}\label{eq:attn-sum}
\bm{u}_j\;\coloneqq\;\sum_{i:(i,j)\in\Nb}\bm{u}_{ij},
\end{equation}
and finally projected through a per-degree channel mixer to the block-output width $C$:
\begin{equation}\label{eq:attn-proj-out}
\Delta\bh^{\mathrm{attn}}_j\;\coloneqq\;L^{\mathrm{ch,out}}(\bm{u}_j)\in\R^{D\times C}.
\end{equation}
The block $L^{\mathrm{ch,out}}$ is again a degree-preserving channel mixer (\texttt{SO3Linear}); it does not couple different $l$'s.

\subsection{\texorpdfstring{$\SO(3)$ equivariance of the attention sublayer}{SO(3) equivariance of the attention sublayer}}\label{sec:attn-equiv}

\begin{proposition}[Equivariance]\label{prop:attn-equiv}
Let $g\in\SO(3)$ and let $\bh\mapsto g\cdot\bh$ be the action of \eqref{eq:feature-action}.
Then $\Delta\bh^{\mathrm{attn}}_j(g\cdot\bh,g\cdot R)=g\cdot\Delta\bh^{\mathrm{attn}}_j(\bh,R)$.
\end{proposition}

\begin{proof}[Sketch]
Under $g$, each edge frame transforms as $\bR_{ij}\mapsto g\bR_{ij}$ (the columns rotate with $g$).
The truncated Wigner $\widetilde\bD_{ij}$ then transforms as $\widetilde\bD_{ij}\mapsto\widetilde\bD_{ij}\,\bD(g)^{-1}$ at the rotation-to-local step and as $\widetilde\bD_{ij}^{-1}\mapsto\bD(g)\,\widetilde\bD_{ij}^{-1}$ at the rotation-back step.
Inside the local frame everything is $\SO(2)_{\by}$-equivariant and the local-frame value $\bm{v}^{\mathrm{loc}}_{ij}$ is therefore invariant under the global $g$.
Attention logits in \eqref{eq:attn-logit} depend only on the $\SO(2)_{\by}$-invariant scalar features $\bm{a}_{ij}$ and so are global-rotation-invariant.
Multiplying the $g$-cancellations and re-rotating in \eqref{eq:attn-back} produces the claimed equivariance.
\end{proof}

% =============================================================
\section{Equivariant Layer Normalisation}\label{sec:eqLN}

EquiformerV3 uses an equivariant normalisation scheme on the per-atom feature.
The default \texttt{merge\_layer\_norm} acts as follows.
For each atom $i$ and each degree $l$, define the per-degree variance
\begin{equation}\label{eq:eqLN-var}
\sigma_{i,l}^2\;\coloneqq\;\frac{1}{(2\min(l,\mmax)+1)\,C}\sum_{m,c}\bigl(\bh_{i,(l,m),c}\bigr)^2,
\end{equation}
where the sum runs over $|m|\le\min(l,\mmax)$ and $c=1,\dots,C$.
Define the normalised feature and apply a per-degree affine gain $\gamma_l\in\R^C$:
\begin{equation}\label{eq:eqLN-norm}
\bh_{i,(l,m),c}\;\leftarrow\;\gamma_{l,c}\,\frac{\bh_{i,(l,m),c}}{\sqrt{\sigma_{i,l}^2+\epsilon}}.
\end{equation}
For the $l=0$ slot, the centring is enabled (subtract the per-channel mean of $\bh_{i,0,c}$ across the atom before normalising).
The variant \texttt{merge\_rms\_norm} omits centring entirely.

% =============================================================
\section{The Transformer Block}\label{sec:block}

A pre-norm transformer block $\Hb^{\mathrm{TB}}_\ell:\R^{N\times D\times C}\to\R^{N\times D\times C}$ is the composition of an attention residual update and an FFN residual update:
\begin{equation}\label{eq:tb-block}
\Hb^{\mathrm{TB}}_\ell\;\coloneqq\;R^{\mathrm{ffn}}_\ell\,\circ\,R^{\mathrm{attn}}_\ell,
\end{equation}
where the two residual updates are defined by
\begin{align}
R^{\mathrm{attn}}_\ell(\bh)\;&\coloneqq\;\bh\;+\;\mathrm{DropPath}\bigl(\Hb^{\mathrm{attn}}_\ell(\LN_1(\bh))\bigr),\label{eq:tb-1}\\
R^{\mathrm{ffn}}_\ell(\bh)\;&\coloneqq\;\bh\;+\;\mathrm{DropPath}\bigl(\Hb^{\mathrm{ffn}}_\ell(\LN_2(\bh))\bigr).\label{eq:tb-2}
\end{align}
Here $\Hb^{\mathrm{attn}}_\ell$ is the equivariant graph-attention sublayer of \S\ref{sec:attn}, $\Hb^{\mathrm{ffn}}_\ell$ is the equivariant feed-forward sublayer of \S\ref{sec:effn}, $\LN_1,\LN_2$ are equivariant layer norms of \S\ref{sec:eqLN}, and $\mathrm{DropPath}$ is the standard stochastic depth applied at the level of whole atoms (one Bernoulli mask per atom per layer, sharing across the irreducible axis).
At inference the dropouts are the identity, so $\Hb^{\mathrm{TB}}_\ell$ is deterministic.

The forward pass through the stack is
\begin{equation}\label{eq:stack}
\bh^{(L)}\;=\;\Hb^{\mathrm{TB}}_L\circ\Hb^{\mathrm{TB}}_{L-1}\circ\cdots\circ\Hb^{\mathrm{TB}}_1\bigl(\bh^{(0)}\bigr),
\end{equation}
followed by a final equivariant layer norm
\begin{equation}\label{eq:stack-out}
\bh^{(L^\star)}\;=\;\LN_{\mathrm{out}}\bigl(\bh^{(L)}\bigr).
\end{equation}

% =============================================================
\section{Energy, Force and Stress Heads}\label{sec:heads}

\subsection{Energy head}\label{sec:energy}

The per-atom energy is read off the $l=0$ slot of the post-norm feature through a two-layer scalar MLP:
\begin{equation}\label{eq:energy-mlp}
\widetilde E_i\;\coloneqq\;\mathrm{Linear}_2\bigl(\SiLU(\mathrm{Linear}_1(\bh^{(L^\star)}_{i,l=0}))\bigr)\in\R,
\end{equation}
with hidden width $C_h^{\mathrm{ffn}}$ (default $128$) and a single scalar output.
The structure-level energy is
\begin{equation}\label{eq:energy-sum}
E_\Theta(Z,R)\;\coloneqq\;\frac{1}{\overline N}\sum_{i=1}^{N}\widetilde E_i,
\end{equation}
where $\overline N$ is a fixed dataset average atom count (\texttt{avg\_num\_nodes}, equal to $77.8$ in the public checkpoint).

\subsection{Direct force head}\label{sec:force-direct}

When direct prediction is selected, EquiformerV3 instantiates an additional equivariant graph-attention head (\S\ref{sec:attn}) producing a single $\SO(3)$-equivariant output channel:
\begin{equation}\label{eq:force-head}
\bm{F}_i^{\mathrm{irr}}\;\coloneqq\;\Hb^{\mathrm{attn,force}}\bigl(\bh^{(L^\star)}\bigr)_i\in\R^{D\times 1}.
\end{equation}
The force vector is read off the $l=1$ slot:
\begin{equation}\label{eq:force-l1}
\bmF_i\;\coloneqq\;\bigl[\bm{F}_i^{\mathrm{irr}}\bigr]_{l=1}\in\R^3.
\end{equation}
This object transforms as a rank-$1$ tensor under $\SO(3)$.
In the default release the force head uses the \texttt{gate} activation (rather than the \texttt{sep-merge\_gates2\_swiglu} of the bulk blocks), which is a simpler $\Sph$-grid variant tailored to single-channel outputs.

\subsection{Gradient force head}\label{sec:force-grad}

When the gradient-prediction path is selected, no separate force head is instantiated; forces are obtained by autograd:
\begin{equation}\label{eq:force-grad}
\bmF_i\;\coloneqq\;-\nabla_{\bR_i}\,E_\Theta(Z,R),
\end{equation}
computed once per forward pass with $\mathtt{create\_graph=self.training}$ to support second-order optimisation.

\subsection{Direct stress head}\label{sec:stress-direct}

When direct stress is selected, an additional equivariant FFN (\texttt{FeedForwardNetworkStressHead}) produces a tensor in $\R^{D\times 1}$, the first $(\lmax+1)^2$ irreducible coefficients of which are passed through a fixed $9\times 9$ change-of-basis matrix $\bm{M}_{\mathrm{stress}}$:
\begin{equation}\label{eq:stress}
\sigma_\Theta(Z,R)\;\coloneqq\;\mathrm{vec}^{-1}\Bigl(\bm{M}_{\mathrm{stress}}\,\bigl[\Hb^{\mathrm{ffn,stress}}(\bh^{(L^\star)})\bigr]_{1:9}\Bigr)\in\R^{3\times 3},
\end{equation}
yielding a symmetric stress tensor in Voigt order.
The matrix $\bm{M}_{\mathrm{stress}}$ is the fixed isometry between the $l=0,1,2$ irreducible blocks and the basis of symmetric $3\times 3$ matrices, given explicitly in \texttt{output\_block.py:get\_stress\_cg\_change\_mat}.

\subsection{Gradient stress head}\label{sec:stress-grad}

When the gradient path is selected for stress, the model differentiates the total energy with respect to a symmetric displacement tensor:
\begin{equation}\label{eq:stress-grad}
\sigma_\Theta(Z,R)\;\coloneqq\;\frac{1}{|\Omega|}\nabla_{\bm{u}}\,E_\Theta(Z(\bm{u}),R(\bm{u}))\Big|_{\bm{u}=0},
\end{equation}
where $\bm{u}\in\R^{3\times 3}$ is a symmetric strain perturbation applied to both atomic positions ($\bR_i\to\bR_i+\bm{u}\bR_i$) and the cell vectors, and $|\Omega|$ is the cell volume.

% =============================================================
\section{Properties}\label{sec:props}

\paragraph{Translation invariance.}
$E_\Theta(Z,R+\bm{t})=E_\Theta(Z,R)$ because only $\br_{ij}=\bR_j-\bR_i$ enters the construction.

\paragraph{Permutation invariance.}
$E_\Theta$ is invariant under $S_N$ because the sum in \eqref{eq:energy-sum} is over atom indices and each per-atom contribution depends only on its neighbourhood through symmetric aggregation operators.

\paragraph{Rotational invariance / equivariance.}
By Proposition~\ref{prop:attn-equiv} and the analogous statement for $\Hb^{\mathrm{ffn}}$ (which follows from $L^{\mathrm{ch}}$ being channel-only and $\bm{S},\bm{S}^{\dagger}$ being independent of the global frame), each transformer block commutes with the global rotation action.
The $l=0$ slot of the per-atom feature is therefore rotation-invariant, so the scalar MLP of \eqref{eq:energy-mlp} produces a rotation-invariant total energy.
The $l=1$ slot of the force head transforms as a vector, so the direct force is $\SO(3)$-equivariant; in the gradient path the equivariance of the force follows from the chain rule applied to a rotationally invariant energy.

\paragraph{Smoothness at the cutoff.}
The envelope $\varphi$ of \eqref{eq:env} satisfies $\varphi(\rcut)=\varphi'(\rcut)=\varphi''(\rcut)=0$.
Because $\varphi$ multiplies both the attention numerator and the radial features at the edge-degree embedding, every contribution from a vanishing edge is $C^2$-smooth in $r_{ij}$.
This is weaker than DPA4's $C^3$ guarantee from the inner-clamp construction but sufficient for the smooth gradient and Hessian of $E_\Theta$ required by direct stress prediction and second-order training.

\paragraph{No spurious dependence on the random secondary axis.}
The random vector $\bu^\star$ enters $\bR_{ij}$ only through the first and third columns, which span the plane perpendicular to $\bu_x=\brh_{ij}$.
Because every in-frame operator is $\SO(2)_{\by}$-equivariant (Remark~\ref{rmk:rand-frame}), the dependence of the output on a rotation of $\bu^\star$ about $\bu_x$ cancels exactly.
The architecture therefore produces deterministic outputs in the ideal-floating-point limit; the random sampling is purely for numerical conditioning of the Euler decomposition.

% =============================================================
\section{Comparison: EquiformerV3 vs.\ DPA4}\label{sec:compare}

We now place EquiformerV3 side-by-side with DPA4 of the companion document.
Throughout, ``DPA4 \S X'' refers to the section of archi\_dpa4.tex with that label; ``E3 \S Y'' refers to this document.
All differences are at the level of design choices within a common framework: both architectures are $\SO(3)$-equivariant, CG-free in the sense of Regime II of DPA4 \S15.1, and operate by a Wigner-rotation to an edge-aligned local frame followed by an $\SO(2)$-equivariant linear--nonlinear--linear sandwich.

\subsection{The common framework}\label{sec:compare-common}

The two architectures share the following ingredients:
\begin{itemize}[leftmargin=2em]
\item Real spherical-harmonic basis for $\SO(3)$ irreducibles, truncated at $\lmax$ and (after rotation) at $\mmax$, with the $m$-major reduced layout.
\item Per-edge $\SO(3)$ frame $\bR_{ij}$ that aligns one Cartesian axis with $\brh_{ij}$, and the corresponding block-diagonal Wigner $\widetilde\bD_{ij}$ used both as the rotation-to-local and (transposed plus a per-$l$ rescale) as the rotation-back operator.
\item Restriction to Regime II of DPA4 \S15.1: messages are equivariant linear maps $\bh_j\mapsto\bm{m}_{ij}$ parameterised by an $\SO(2)$-equivariant kernel $L^{\mathrm{SO2}}$ surrounded by $\widetilde\bD_{ij}$, with no Clebsch--Gordan tensor product of two equivariant features.
\item Per-edge $\SO(2)$-equivariant linear operators that mix different $l$'s within each $|m|$-stratum (the cross-$l$ mixing of Regime II is realised in both architectures).
\item A smooth radial cutoff envelope.
\item Equivariance achieved by separating linear operators (constrained) from nonlinearities (applied either pointwise on scalars or on a $\Sph$-grid, both invariant under $\SO(2)_{\by}$).
\end{itemize}

\subsection{Differences, component by component}\label{sec:compare-diff}

\paragraph{Frame axis.}
DPA4 (``Local Frame on a Directed Edge'' in archi\_dpa4.tex) aligns $\brh_{ij}$ with the local $\bz$-axis and writes the in-frame symmetry group as $\SO(2)_{\bz}$; the rotation operator is constructed by a two-chart NLERP blend on $\Sph$ that produces a deterministic, $C^3$-smooth quaternion as a function of $\brh_{ij}$.
EquiformerV3 (E3 \S\ref{sec:edgerot}) aligns $\brh_{ij}$ with the local $\by$-axis and assembles the rotation matrix from a freshly drawn random perpendicular per edge.
The randomness is harmless under the $\SO(2)_{\by}$-equivariance of every in-frame operator (Remark~\ref{rmk:rand-frame}), but only $C^0$-smooth as a function of $\brh_{ij}$ (the random direction is discontinuous in $\brh_{ij}$).
DPA4's frame is $C^3$ in $\brh_{ij}$; EquiformerV3's frame is not.
For the architectures' actual outputs the smoothness gap is closed by $\SO(2)$-equivariance: any $\SO(2)_{\by}$-invariant function of $\bm{x}^{\mathrm{loc}}_{ij}$ is automatically continuous in $\bu^\star$ (it does not see $\bu^\star$ at all).

\paragraph{Wigner construction.}
DPA4 builds $\widetilde\bD_{ij}$ directly from the unit quaternion via the closed-form real-basis Wigner-$D$ formula.
EquiformerV3 builds it from ZYZ Euler angles via $e^{-i\alpha M_z}d(\beta)e^{-i\gamma M_z}$ in the real basis, with $d(\beta)$ contracted against a precomputed Jucys--Murphy tensor stored on disk.
The two constructions agree on the resulting $\widetilde\bD_{ij}$ up to numerical noise.

\paragraph{Radial basis.}
DPA4 uses learned smooth-bump radial functions with explicit $C^3$ continuity at the inner clamp and the outer cutoff.
EquiformerV3 uses $K_r=600$ Gaussian basis functions \eqref{eq:gauss} with a single fixed width, followed by an MLP \eqref{eq:radfn}; this is a strictly richer family at finite width but loses the analytic smoothness control.
The two architectures fix their radial mixing weights at different points: DPA4 ties weights across $m$ via the per-degree rescale $\rho_l$, while EquiformerV3 toggles this with the \texttt{use\_rad\_l\_parametrization} flag (default: per-$l$, matching DPA4's tying).

\paragraph{Cutoff envelope.}
DPA4 uses a septic Hermite envelope $s_p(r)$ that is $C^3$ at both ends of the cutoff region.
EquiformerV3 uses a degree-$p+2$ polynomial of \eqref{eq:env} that is $C^2$ at the outer cutoff and unbounded near $r=0$ (no inner clamp).
For molecular configurations with bond lengths far from any divergence the two are interchangeable in practice; for pathological short-distance configurations DPA4's inner-clamp and ZBL bridging confer extra stability.

\paragraph{ZBL short-range bridging.}
DPA4 includes an additive analytic ZBL pair potential bridged smoothly into the learned energy through the source-freeze propagation gate $\eta_j$ (``Source-Freeze Gate'' and ``ZBL Pair Potential'' in archi\_dpa4.tex).
EquiformerV3 has no such bridging; the energy is purely the learned scalar of \eqref{eq:energy-sum}.
Consequently EquiformerV3 has no guarantee of correct short-range physics outside the support of its training distribution.

\paragraph{Initial embedding of higher-$l$ slots.}
Both architectures initialise the per-atom feature with a species embedding in the $l=0$ slot and zeros elsewhere, then add a one-shot edge-degree embedding to populate the $l>0$ slots.
DPA4 calls this the Geometric Initial Embedding (GIE).
The two GIEs differ only in the radial parameterisation (DPA4's bump basis vs.\ EquiformerV3's Gaussian basis) and in the aggregation normaliser ($1/\sqrt{\overline d}$ in DPA4, $1/\overline d$ in EquiformerV3 \eqref{eq:edgedeg-agg}).

\paragraph{Edge message: attention vs.\ pure convolution.}
This is the principal architectural divergence.
EquiformerV3 (E3 \S\ref{sec:attn}) computes a softmax attention weight $\alpha^{(h)}_{ij}\in[0,1]$ per head from an $\SO(2)_{\by}$-invariant scalar logit and applies it multiplicatively to the local-frame value; the per-target aggregation is a sum of attention-weighted local-frame messages.
DPA4's edge-convolution block (``Edge Convolution in the Local Frame'' in archi\_dpa4.tex) does \emph{not} use softmax attention; aggregation is a scalar-keyed convex combination $\mathcal{A}$ of per-edge messages with mixing weights driven by a softmax over the depth-attention channel, plus a residual.
The DPA4 attention is depth-attention over feature streams within the same block; EquiformerV3 attention is graph-attention over edges incident to the same target.
Both are global-rotation invariant in the weighting; both are smooth in the edge geometry through the envelope, but only EquiformerV3 includes an explicit per-edge attention softmax \eqref{eq:attn-softmax}.

\paragraph{Source--target asymmetry in the edge message.}
EquiformerV3 mixes both source and target node features into the per-edge message via the concat (or add-merge) operation of \eqref{eq:attn-edge}.
DPA4 uses only the source feature $\bh_j$ in the per-edge message, with the target's role entering only through species embeddings in the radial.
This is a strict architectural asymmetry: EquiformerV3 messages are functions of $(\bh_i,\bh_j,r_{ij},Z_i,Z_j)$; DPA4 messages are functions of $(\bh_j,r_{ij},Z_i,Z_j)$ only.

\paragraph{$\SO(2)$ linear operator.}
The two architectures use the \emph{same} $\SO(2)$-equivariant linear operator class on the truncated $m$-major reduced layout: an unconstrained dense block on the $m=0$ stratum, plus the real form of complex multiplication $\bigl[\begin{smallmatrix}\bU_m&-\bV_m\\\bV_m&\bU_m\end{smallmatrix}\bigr]$ on each $|m|>0$ stratum with cross-$l$ mixing allowed within the stratum.
Both rescale the $|m|>0$ weights by $1/\sqrt{2}$ to compensate for the doubled real parameter count.
Remark~\ref{rmk:so2-equiv} gives the precise statement, including the few incidental differences (multi-stream focus axis, side-channel scalars, bias-correction tied to the radial--envelope product) that do not alter the equivariance class.

\paragraph{In-frame nonlinearity.}
Both architectures implement the same $\Sph$-grid sandwich for the FFN nonlinearity, and DPA4 explicitly enumerates three FFN variants in archi\_dpa4.tex \S11 (``The Equivariant Feed-Forward Network''):
\emph{Variant 1} (gated FFN, eq.~(58) of archi\_dpa4.tex) -- the default \texttt{GatedActivation} path of \texttt{sezm\_nn/ffn.py:EquivariantFFN}, which gates $l>0$ slots by a sigmoid of $l=0$ features;
\emph{Variant 2} (spherical-grid FFN, eq.~(59) of archi\_dpa4.tex) -- the \texttt{grid\_mlp=True} path, which up-projects per-degree, sends the higher-$l$ part through the grid sandwich $\mathsf{P}^{\dagger}\circ\Psi\circ\mathsf{P}$ (with $\Psi$ a per-grid-point channel-mixing MLP), adds a separate scalar SwiGLU branch on the $l=0$ slot, then down-projects per-degree;
\emph{Variant 3} (spherical-grid SwiGLU activation, eq.~(60) of archi\_dpa4.tex) -- the \texttt{s2\_activation=True} path, where the activation itself is a grid-SwiGLU built from \texttt{SwiGLUS2Activation}, splitting the hidden width into value and gate, evaluating the gate on $\Sph$, multiplying pointwise, and projecting back.
EquiformerV3's FFN \eqref{eq:effn-1} is structurally Variant~2: the same per-degree up-/down-projection, the same e3nn \texttt{ToS2Grid}/\texttt{FromS2Grid} construction for $\mathsf{P},\mathsf{P}^{\dagger}$, the same per-degree truncation rescale for $\mmax<\lmax$, the same per-grid-point pointwise MLP, and the same scalar-merge to the $l=0$ slot.
The differences from DPA4's Variant~2 are: (i) DPA4 subtracts the $l=0$ slice before projecting to the grid ($\mathsf{P}(\bh-\bh|_{l=0})$ in eq.~(59) of archi\_dpa4.tex), so the grid path acts only on $l\ge 1$ and the $l=0$ slot is processed disjointly by the scalar SwiGLU branch, whereas EquiformerV3 keeps $l=0$ inside the grid input and adds the scalar branch on top of it;
(ii) EquiformerV3's grid MLP multiplies the SwiGLU output by a sigmoid scalar gate $\Sigmoid(\bm{w}_g\bm{s})$ derived from $\bh_{l=0}$ \emph{inside} the grid (eq.~\eqref{eq:gridmlp}), while DPA4's \texttt{PointwiseGridMLP} has no scalar gate inside the grid, the scalar information entering only through the additive merge;
(iii) the grid path is opt-in (\texttt{grid\_mlp=False} default) in DPA4 and the default in EquiformerV3;
(iv) the latitude grid in EquiformerV3 is dense ($N_\beta=20$) and the longitude grid sparse ($N_\alpha=8$ in attention, $N_\alpha=20$ in FFN), while DPA4's grid resolution choices are smaller in both axes.
For the attention sublayer's in-frame activation, EquiformerV3's \texttt{sep-merge\_gates2\_swiglu} \eqref{eq:attn-act} corresponds to DPA4's Variant~3 used as the inter-mixer activation inside the \texttt{SO2Convolution} stack of $\SO(2)$-linear layers.

\paragraph{Layer normalisation.}
DPA4 uses a per-degree equivariant RMS-norm without centring; EquiformerV3 defaults to the merge layer-norm of \S\ref{sec:eqLN} with centring on $l=0$ and RMS-style normalisation on $l>0$.
The variants \texttt{merge\_rms\_norm} and \texttt{sep\_layer\_norm} bring EquiformerV3 closer to DPA4's choice.

\paragraph{Transformer block layout.}
EquiformerV3 uses a pre-norm attention $+$ pre-norm FFN block \eqref{eq:tb-1}--\eqref{eq:tb-2} with stochastic depth (DropPath) applied at the level of whole atoms.
DPA4's interaction block has a sandwich-norm structure (norm$\to$op$\to$norm$\to$residual) with no stochastic depth, and the depth-attention residual aggregation $\mathcal{A}$ replaces the additive residual.

\paragraph{Energy normaliser.}
EquiformerV3 divides the summed per-atom energy by $\overline N$ (the dataset's average atom count); DPA4 does not perform this division, instead relying on a learned bias and the consistent per-atom scale of the readout.

\paragraph{Forces and stress.}
EquiformerV3 supports both a direct prediction head (additional equivariant attention/FFN producing $l=1$ vectors and stress-basis coefficients) and a gradient prediction by autograd; the two paths are mutually exclusive at instantiation.
DPA4 produces forces and stress only by autograd through the energy, and is consequently always conservative.
The direct heads of EquiformerV3 are strictly more general (they can produce non-conservative force fields) but lose the conservation guarantee.

\paragraph{Maximum degrees and orders.}
The released checkpoints set $\lmax=6,\mmax=2$ for EquiformerV3; DPA4 typically operates at smaller $\lmax$ and $\mmax$ with a tighter $C$.
Because in-frame compute scales as $\mathcal{O}(\lmax\cdot\mmax\cdot C\cdot C_h)$ in Regime II, both architectures land in the same regime asymptotically.

\subsection{Operator-level identification}\label{sec:compare-op}

In the notation of DPA4 \S15.1, both architectures fit a common per-edge message template
\begin{equation}\label{eq:common-msg}
\bm{m}_{ij}\;=\;K\;\widetilde\bD_{ij}^{-1}\;\Hb^{\mathrm{loc}}\bigl(\widetilde\bD_{ij}\,\bigl(\widetilde\brho_{ij}\odot\bh^{\mathrm{src/tgt}}_{ij}\bigr)\bigr),
\end{equation}
where $\widetilde\brho_{ij}\in\R^{(\lmax+1)\times \cdot}$ is the edge-conditioned per-$l$ radial scaling, $K$ is a scalar normaliser, and $\Hb^{\mathrm{loc}}$ is an $\SO(2)$-equivariant nonlinear operator.
The two architectures specialise this template as follows:

\medskip
\begin{tabular}{p{4.5cm}|p{5cm}|p{5cm}}
\hline
\textbf{Slot} & \textbf{DPA4} & \textbf{EquiformerV3}\\
\hline
$\bh^{\mathrm{src/tgt}}_{ij}$ & $\bh_j$ (source only) & $[\bh_i\,\|\,\bh_j]$ (concat) or $\widetilde\brho^{\mathrm{src}}\bh_i+\widetilde\brho^{\mathrm{tgt}}\bh_j$ (add-merge) \\
$\widetilde\brho_{ij}$ & per-$(l,c)$ bump basis & per-$(l,c)$ Gaussian-MLP \\
Local frame axis & $\bz$ via $C^3$ NLERP & $\by$ via random perpendicular \\
$\Hb^{\mathrm{loc}}$ & gated SwiGLU sandwich, no attention & SO(2)-linear $\to$ SwiGLU on $\Sph$-grid $\to$ SO(2)-linear, multiplied by softmax attention weight $\alpha^{(h)}_{ij}$\\
$K$ & per-$l$ rescale $\rho_l$ \eqref{eq:wignerD-trunc-inv} & per-$l$ rescale $s_l$ \eqref{eq:wignerD-trunc-inv}\\
Aggregation & scalar-keyed convex combination $\mathcal{A}$ over depth $+$ smooth envelope $s_p$ & sum-over-edges $+$ envelope-rescaled softmax $+$ DropPath \\
Short-range bridging & ZBL $+$ source-freeze gate $\eta_j$ & none \\
Force prediction & autograd & autograd or separate $l=1$ head \\
\hline
\end{tabular}

\medskip
Both the Regime-II message of archi\_dpa4.tex and the per-edge message of EquiformerV3 \eqref{eq:common-msg} occupy the same expressiveness slot (Regime II): direction-dependent, cross-$l$-mixing, equivariant-linear maps on a single equivariant feature.
The architectural divergences enumerated above are concerned with the parameterisation of $\widetilde\brho_{ij}$, $L^{\mathrm{SO2}}$, $\Hb^{\mathrm{loc}}$, and the aggregation, not with the expressiveness class.

\subsection{Where the architectures genuinely differ}\label{sec:compare-genuine}

Three of the differences above are not parameterisation choices but architectural commitments:
\begin{enumerate}[leftmargin=2em]
\item \emph{Smoothness of the edge frame.}
DPA4's frame is $C^3$ in $\brh_{ij}$; EquiformerV3's is $C^0$ but $\SO(2)$-invariantly $C^\infty$.
For applications requiring smooth force gradients (Hessians, free-energy second derivatives), the DPA4 frame is preferable.
\item \emph{Short-range bridging.}
DPA4's ZBL bridging is a non-perturbative addition to the energy that guarantees correct physics in the short-range limit; EquiformerV3 has no such mechanism.
\item \emph{Per-edge attention.}
EquiformerV3 uses softmax attention over edges; DPA4 uses an unweighted sum (with smooth envelope) over edges and a separate depth-attention residual.
The qualitative effect of the EquiformerV3 attention is to make the per-edge contribution adaptive to the neighbourhood; the qualitative effect of the DPA4 depth-attention is to make the per-feature-stream contribution adaptive to the depth.
The two attention mechanisms are not interchangeable in the literal sense, although both can in principle be subsumed by a sufficiently expressive Regime-II kernel of archi\_dpa4.tex.
\end{enumerate}

\subsection{What is shared}\label{sec:compare-shared}

All of the following are common to both architectures:
\begin{itemize}[leftmargin=2em]
\item The decision to live in DPA4 Regime II rather than Regime III (no Clebsch--Gordan coefficients explicitly invoked).
\item The use of a per-edge Wigner-$D$ rotation as the carrier of angular information.
\item Per-degree channel mixers as the only $\SO(3)$-linear operator on global-frame features (archi\_dpa4.tex's $L^{\mathrm{ch}},L^{\mathrm{deg}}$ coincide with EquiformerV3's \texttt{SO3Linear}).
\item Real-basis spherical harmonics throughout.
\item Energy by per-atom $l=0$ MLP plus structure sum.
\item Smooth radial cutoff at $r=\rcut$.
\end{itemize}

\subsection{Summary}\label{sec:compare-summary}

EquiformerV3 and DPA4 are both Regime-II CG-free architectures whose per-edge messages live in the same equivariant operator class.
They differ on the parameterisation of that operator (Gaussian-MLP vs.\ bump basis), on the construction of the edge frame ($C^0$ random vs.\ $C^3$ NLERP), on the use of attention (graph-attention over edges vs.\ depth-attention over features), on the inclusion of short-range physics (none vs.\ ZBL bridging), and on the output head structure (direct + gradient vs.\ gradient only).
None of these differences alter the asymptotic expressiveness class; each is a deliberate trade-off between smoothness, expressive headroom, training stability, and end-application demands.

% =============================================================
\section*{Notation Index}

\begin{tabular}{ll}
$N$ & number of atoms\\
$Z_i\in\{1,\dots,Z_{\max}\}$ & atomic species of atom $i$\\
$\bR_i\in\R^3$ & position of atom $i$\\
$\br_{ij}=\bR_j-\bR_i$, $r_{ij}=\|\br_{ij}\|$, $\brh_{ij}=\br_{ij}/r_{ij}$ & edge vector, distance, direction\\
$\rcut$ & outer cutoff radius\\
$\lmax,\mmax\in\Z_{\geq 0}$ & maximum degree, maximum truncation order\\
$D=(\lmax+1)^2$, $D_m=\sum_l(2\min(l,\mmax)+1)$ & full and truncated dimensions\\
$\iota(l,m)=l^2+l+m$ & linear packing of $(l,m)$\\
$C$ & feature channel width\\
$H$ & number of attention heads\\
$C^{\mathrm{a}},C^{\mathrm{v}}$ & per-head alpha- and value-channel widths\\
$C_h$ & MLP hidden width\\
$L$ & number of transformer blocks\\
$N_\beta,N_\alpha$ & grid latitude/longitude resolutions\\
$\bR_{ij}\in\SO(3)$ & edge-aligned rotation matrix \eqref{eq:edgerot}\\
$\bD_{ij}\in\R^{D\times D}$ & block-diagonal Wigner-$D$ at $\bR_{ij}$\\
$\widetilde\bD_{ij}\in\R^{D_m\times D}$ & truncated $m$-primary Wigner \eqref{eq:wignerD-trunc}\\
$\widetilde\bD_{ij}^{-1}\in\R^{D\times D_m}$ & rescaled inverse Wigner \eqref{eq:wignerD-trunc-inv}\\
$\varphi(r)$ & polynomial envelope \eqref{eq:env}\\
$\eta(r)$ & Gaussian smearing \eqref{eq:gauss}\\
$\Hb_{\mathrm{rad}}$ & radial MLP \eqref{eq:radfn}\\
$L^{\mathrm{ch}}$ & per-degree channel mixer (\texttt{SO3Linear})\\
$L^{\mathrm{SO2}}=(L_0,L_1,\dots,L_{\mmax})$ & $\SO(2)_{\by}$-equivariant linear operator \eqref{eq:LSO2}\\
$\bm{S},\bm{S}^{\dagger}$ & grid evaluation and pull-back \eqref{eq:sample}--\eqref{eq:grid-pull}\\
$\mathcal{M}_{\Sph}$ & grid MLP \eqref{eq:gridmlp}\\
$\Hb^{\mathrm{attn}}_\ell,\Hb^{\mathrm{ffn}}_\ell$ & block sublayers\\
$\Hb^{\mathrm{TB}}_\ell$ & one transformer block \eqref{eq:tb-1}--\eqref{eq:tb-2}\\
$\alpha^{(h)}_{ij}$ & per-edge per-head attention weight \eqref{eq:attn-softmax}\\
$\overline N,\overline d$ & dataset-derived atom-count and degree averages\\
$\bmF_i,\sigma_\Theta$ & per-atom force and stress\\
\end{tabular}

\end{document}
